

\section{Apprentissages et retours d'expérience}

Ce projet nous a permis de développer bien plus que des compétences techniques. Il a constitué une expérience pédagogique complète et immersive, qui a fait appel à notre autonomie, notre esprit d’équipe, notre rigueur et notre capacité à résoudre des problèmes concrets dans un cadre structuré.



\subsection{Perspectives personnelles}

Ce projet a été l’occasion de donner du sens aux notions abordées en cours, souvent jugées abstraites. Voir des agents se déplacer, réagir, coopérer ou poursuivre, a rendu nos algorithmes vivants et nos classes Java tangibles.

Il nous a permis de mieux comprendre la logique des grands systèmes logiciels et de nous préparer aux défis à venir dans notre parcours universitaire et professionnel. Il nous a aussi sensibilisés aux enjeux d’organisation, de lisibilité du code, et de communication interdisciplinaire dans une équipe.

En résumé, cette expérience a été extrêmement formatrice et motivante. Elle nous a permis de mesurer le chemin parcouru depuis le début de notre formation, tout en nous donnant envie de pousser encore plus loin nos compétences dans les domaines du génie logiciel, ou des simulations interactives.